\section{Introduction}
\label{sec:introduction}

Choosing a forecasting architecture for solar irradiance requires evidence
across the temporal resolutions and horizons at which it will be used.
This article compares TimesNet and DilatedRNN to identify a preferred
implementation and the conditions that support that choice.

Pathways consistent with deep emissions reductions combine a rapid expansion
of low-emission electricity with the electrification of end uses
\cite{ipcc2022energy,irena2024weto}. Declining renewable-energy costs support
this transition \cite{luderer2022}, while scenario-based evidence suggests
that solar photovoltaics may develop substantial self-sustaining growth
before mid-century \cite{nijsse2023}. As deployment grows, the
weather-dependent and diurnal variability of solar energy makes forecasts
across different time horizons increasingly important for planning,
scheduling, and energy management.

Industrial decarbonization requires combinations of energy efficiency,
material changes, low-carbon fuels, electrification, and low-emission
electricity \cite{ipcc2022industry,rissman2020}. Direct electrification can
reduce emissions from suitable industrial heat uses, but its potential
depends on process temperature, infrastructure, and the carbon intensity of
the electricity supply \cite{madeddu2020}. On-site and procured renewable
electricity can therefore support manufacturing alongside other energy
options \cite{irena2014manufacturing,padmanathan2018}, although adoption
faces site-specific financial, regulatory, technical, and organizational
barriers \cite{dhingra2023}.

Policy and corporate disclosures illustrate this industrial context. In the
European Union, Directive (EU) 2024/1275 establishes a phased framework for
solar installations on specified building categories when technically,
economically, and functionally suitable \cite{eu2024epbd}; it does not apply
to the Brazilian, Mexican, or United States locations examined here.
Nestl\'e reports that its KitKat production lines in Ca\c{c}apava use
predominantly solar electricity \cite{nestle2024biometano}; Mercedes-Benz
describes rooftop photovoltaics and storage at Factory~56
\cite{mercedesfactory56}; and Tesla and Ford report solar installations at
selected manufacturing sites \cite{tesla2021impact,ford2023valencia}.

Solar forecasting serves different operational needs. Intra-hour forecasts
can track cloud motion, hourly and day-ahead forecasts support short-term
scheduling, and daily or monthly aggregates inform longer-term planning
\cite{diagne2013,inman2013}. Both temporal resolution and forecast horizon
must therefore accompany terms such as short-term and long-term
\cite{antonanzas2016,yang2018}. Accuracy for a direct 72-hour sequence of
hourly global horizontal irradiance (GHI) does not establish accuracy for
daily means or a one-step monthly target. Likewise, comparing errors across
hourly and monthly series requires accounting for differences in targets,
sample sizes, and variability.

This article studies GHI at ten geographic points associated with
electric-vehicle manufacturing facilities. The points span different
latitudes and atmospheric regimes, providing a geographically diverse test
bed. Their selection does not establish whether photovoltaic systems are
installed at those facilities. The irradiance records are satellite-derived,
model-based estimates from the National Solar Radiation Database (NSRDB)
\cite{sengupta2018}, with no factory sensor measurements or proprietary
operational data. The forecasts thus characterize the solar resource at the
selected grid cells, rather than factory load or photovoltaic power output.

Cabral et al.\ \cite{cabral2026} evaluated very-long-term GHI forecasting at
40 climatically diverse locations using a GHI-only GPT-Neo pipeline and a
broad set of competitors. Here, the focus is on comparing architectures
across temporal resolutions. Building on the present authors' preliminary studies of
a compact TimesNet for direct hourly forecasting and a dilated recurrent
neural network (DilatedRNN) for one-step monthly forecasting, this article
evaluates both architectures on matched samples within each task.
The preliminary studies provide methodological background; their numerical
results are excluded from the matched comparison. Cabral et al.'s numerical
results are not reused.

The research question is: \emph{which of TimesNet and DilatedRNN is better
supported by a matched multi-resolution evaluation, and at which horizons
and locations does the preference change?} The comparison uses the
cumulative macro-MAE over each task's full output horizon to summarize
task-level outcomes, alongside intermediate horizons, site-level
differences, and seed variability. The four tasks span three resolutions;
their win count is a descriptive summary rather than a pooled error score.

The article makes four contributions:
\begin{enumerate}
  \item it defines a common multi-resolution formulation that distinguishes
  hourly multi-step, daily-mean multi-step, one-step monthly, and exploratory
  six-month forecasting;
  \item it specifies matched samples, chronological partitions, forecast
  origins, direct outputs, and post-processing for the two architectures;
  \item it identifies DilatedRNN as the lower-macro-MAE architecture in
  three of the four full-horizon tasks, with its largest relative advantage
  in one-month forecasting, and examines this preference across sites and
  seeds while retaining the hourly advantage of TimesNet; and
  \item it combines horizon profiles, parameter counts, and reproducibly
  selected forecast cases to characterize the scope of the model choice.
  Climate and independent meteorological records provide descriptive
  context for the cases.
\end{enumerate}
The complete reference-model benchmark is supplied in the Supplementary
Material to place this two-architecture comparison in context.

Section~\ref{sec:related_work} reviews the forecasting literature and
identifies the comparison gap. Section~\ref{sec:background} defines GHI,
temporal aggregation, forecasting horizons, TimesNet, and DilatedRNN.
The following sections describe the data and feature engineering, common
methodology, experimental design, and evaluation metrics, followed by the
results, limitations, and conclusions.
